\documentclass{article}
\usepackage{pathfinder-readable}

\title{Token charges and delayed trade in an LLM double auction}

\begin{document}
\maketitle

\begin{abstract}
This note connects a theory of rewards for AI agents to a proposal to charge traders for the tokens
they generate in a double auction. The thread looked for a test of whether a token charge changes
choices as an actual price would, and for a reason it might fail to bring forward a useful trade. It
found a necessary token-use restriction at a fixed decision and a conditional example in which an
own-token charge leaves a socially costly delay. A reward based on the whole trading path fixes that
example if prices and token counts can be verified. The thread paused because the auction has not
been run with these incentives, and the example's timing and funding assumptions remain untested.
\end{abstract}

\section{Introduction}

Bergemann, Koh and Morris study rewards for AI agents with uncertain aims and abilities \cite{Q}.
They ask whether rewards make agents follow an intended course. A complete sequence of actions can
count as one choice, and one agent's reward can depend on what several agents do. The paper gives no
result for a sequential double auction.

The local auction paper studies the cost of price discovery when language-model buyers and sellers
make quotes \cite{P}. A buyer can accept a standing ask, or bid and wait. Generated tokens cost
resources, so an earlier trade may save later calls. The paper proposes a charge for each trader's
own tokens and asks whether it changes crossing, participation and net welfare. It reports no run
with that charge. Its internal labels Q and P name other source papers; here they mean the assigned
pair.

The connection was to treat a trader's full auction path as an action and its token charge as a
reward deducted from its payoff. This made two questions precise: whether observed choices would
respond to a real charge in the predicted way, and whether that charge would account for tokens
spent by the other trader.

\section{What was done}

The peers first caught the mismatch between the assigned pair and the auction paper's own Q and P
labels. They then considered using the theory paper's peer-scoring method to elicit traders' values,
while mapping a trader's auction path and token charge onto the reward framework. That map supplies
a way to state incentives, not an auction result. They explored a delayed-trade example because the
auction paper already says that an early crossing may save a counterparty's later call. One early
version allowed uncertain later trade and extra calls. It exposed a possible gap between a buyer's
preferred price and the total cost of reaching a match, but depended on unmeasured continuation
probabilities.

The peer-scoring route required an added assumption that values were randomly permuted among
traders. A proposed numerical separation applied to one peer's value, whereas the theory paper's
score uses the whole profile of peers. Direct reports would also replace the standing-book auction
under study. The peers set this route aside.

The next step was a simpler diagnostic. One peer derived a restriction on token demand at two actual
charge rates while holding the decision fixed; the other checked the algebra. The peers then made
the delayed-trade example sharper by requiring a later seller call that the scheduler cannot avoid
and by making its token cost the same whether the seller accepts or declines. Finally, they checked
a reward tied to the complete path that removes both the buyer's reason to wait for a better price
and the cost imposed on the seller. These were derivations and checks, not auction runs.

\section{Results}

For the diagnostic, fix a trader, its value, the book, the other traders' behaviour and the choices
available. Let \(a\) be a choice, \(U(a)\) its expected payoff before the charge, and \(T(a)\) its
expected metered tokens. Suppose these quantities stay the same when the token price changes. If
\(a_0\) is optimal at price \(\tau_0\), and \(a_1\) at a higher price \(\tau_1\), each choice must
beat the other at its own price. Adding those two payoff comparisons gives
\[
(\tau_1-\tau_0)\bigl(T(a_0)-T(a_1)\bigr)\geq 0.
\]
Thus \(T(a_1)\leq T(a_0)\): under these fixed conditions, a higher actual charge cannot make an
optimal trader choose more expected tokens. The same reasoning covers lotteries over choices, using
their expected payoffs and token counts. Emmy derived this restriction and Ada checked it in the
peer notes. It is a necessary test of the additive-charge model, not a prediction for total tokens
or crossing in a live auction, where the book and other traders can change.

The delayed-trade example shows a different limit. Let a buyer value an item at \(v\), let a
seller's cost be \(c\), and let the standing ask be \(p\). The buyer can accept \(p\) now or bid
\(p-\delta\), where \(\delta>0\) is the price cut and \(v>p>p-\delta>c\). Suppose the buyer's two
replies use the same number of tokens. After a bid, the scheduler requires a seller call using
\(K_S>0\) extra tokens. The seller pays for that call whether it accepts or declines, and no further
buyer call is needed. Once called, the seller accepts the bid because its price exceeds \(c\).

Both paths make the same match. At every own-token price \(\tau\geq0\), the buyer prefers to bid: it
saves \(\delta\) on the item without changing its own token bill. Let \(\lambda>0\) convert token
cost into the units of trading surplus. The gross gain from trade is \(v-c\) on either path, but
crossing now raises net welfare by \(\lambda K_S\) because it avoids the seller's call. Ada stated
the example, and Emmy checked the continuation conditions. It is a possible local failure of an
own-token charge, not a finding about traders in the proposed auction.

The peers then checked a repair for this fixed match. For a complete trading path \(a\), let
\(p(a)\) be the execution price and \(C_B(a)\) and \(C_S(a)\) be the buyer's and seller's verifiable
token totals. Give the buyer the additional reward
\[
R_B(a)=p(a)-\lambda C_S(a)-(\lambda-\tau)C_B(a).
\]
The buyer's payoff after its own-token charge and this reward is \(v-\lambda[C_B(a)+C_S(a)]\). It
therefore favours the lower-token path when the match is fixed. The price term matters: charging for
both traders' tokens alone can still leave waiting attractive when \(\delta>\lambda K_S\). The
reward calculation was checked by both peers. It establishes only this local comparison.

\section{What did not hold}

The peer-scoring idea did not justify a result about this auction. It needed a value-assignment rule
absent from the auction paper and a different score for the single-peer calculation. The peers kept
it only as a conditional benchmark.

The token-price restriction needs an actual additive charge and a stable decision problem. An
announced charge may change choices without making the trader maximise the stated payoff. More
tokens at a higher price would reject the combined assumptions without identifying which failed.
Fewer tokens would not establish auction efficiency.

The delayed-trade example needs equal-token buyer replies and a mandatory seller call whose cost is
sunk before the seller chooses. Otherwise a high charge might prevent the later trade. The repair
needs verifiable prices and token counts and might need outside funds. It does not settle private
values, changing matches, seller outside options or the full auction's equilibrium. The theory
paper's coupled-reward example concerns another game. The auction paper already mentions
counterparty savings, so this example sharpens that concern without establishing an empirical
effect.

A targeted search in the ledger found adjacent work on token-budget prompting \cite{tokenbudget} and
language-model preference tests \cite{preferences}. The search was limited and supports no broad
claim of novelty.

\section{Why the thread stopped}

The verifier accepted the fixed-decision inequality and conditional crossing example, but paused
because no charged traders have been tested under the actual scheduler. The reward needs a fixed
match and observable path data; its funding and enforcement are unresolved. The smallest restart is
a metered fixed-book test that varies a real token charge while holding options and other traders'
behaviour fixed. A later paired crossing test could check the seller continuation and compare the
own-token charge with a specified path reward. Full-auction welfare remains a later question.

\bibliographystyle{plainurl}
\bibliography{references}

\end{document}
